% !TEX root = ../LosningsforslagTRES0312.tex
\chapter{Hva er fysikk?}
\label{LF:komp-Chapter:hvaerfysikk}
Dette kapittelet inneholder løsningsforslag til oppgavene i Kapittel~\ref{komp-Chapter:hvaerfysikk}, \textit{Hva er fysikk?}, i kompendiet.

\oppgaver{SI-enheter}{si-enheter}{%
%
Oppgi SI-enheten og enhetssymbolet for følgende fysiske størrelser:
\begin{enumerate}[a)]
    \item Masse
    \item Tid
    \item Temperatur
    \item Elektrisk strøm
    \item Lengde
\end{enumerate}
%
}

\svar[title={Løsningsforslag}]{%
Hver fysisk størrelse har én bestemt SI-enhet, se Tabell~\ref{komp-tab:SI_enheter} i kompendiet.
\begin{enumerate}[a)]
    \item Masse: kilogram, symbol $\mathrm{kg}$.
    \item Tid: sekund, symbol $\mathrm{s}$.
    \item Temperatur: kelvin, symbol $\mathrm{K}$.
    \item Elektrisk strøm: ampere, symbol $\mathrm{A}$.
    \item Lengde: meter, symbol $\mathrm{m}$.
\end{enumerate}
}

\oppgaver{Enhetsregning}{enhetsregning}{%
%
En bil kjører med fart $v = \SI{90}{\kilo\metre\per\hour}$.
\begin{enumerate}[a)]
    \item Gjør om farten til m/s. \textit{Hint: } $\SI{1}{\kilo\metre} = 10^3\,\si{\metre}$ og $\SI{1}{\hour} = \SI{3600}{\second}$.
    \item Kjøreturen varer i $t = \SI{2.5}{\hour}$. Finn tilbakelagt strekning i km og i m.
\end{enumerate}
%
}

\svar[title={Løsningsforslag}]{%
\noindent
\textbf{a)} Vi regner om ved å bruke at $\SI{1}{\kilo\metre}=10^3\,\si{\metre}$ og $\SI{1}{\hour}=\SI{3600}{\second}$:
\[
    v = \SI{90}{\kilo\metre\per\hour} = \frac{90\cdot 10^3\,\si{\metre}}{3600\,\si{\second}} = \doubleunderline{\SI{25}{\metre\per\second}}
\]

\smallskip\noindent
\textbf{b)} Tilbakelagt strekning er fart multiplisert med tid:
\[
    s = v\cdot t = \SI{90}{\kilo\metre\per\hour}\cdot\SI{2.5}{\hour} = \SI{225}{\kilo\metre} \approx \doubleunderline{\SI{2.3e2}{\kilo\metre}}
\]
I meter blir dette
\[
    s = \SI{225}{\kilo\metre}\cdot 10^3\,\frac{\si{\metre}}{\si{\kilo\metre}} \approx \doubleunderline{2.3\cdot 10^5\,\si{\metre}}.
\]
}

\oppgaver{Enhetsregning}{enhetsregning-mane}{%
%
Tyngdeakselerasjonen på månen er $\mathrm{g}_\textrm{M} = \SI{1.62}{\metre\per\second\squared}$. Gjør om til $\si{\kilo\metre\per\minute\squared}$.
%
}

\svar[title={Løsningsforslag}]{%
Vi skal regne om fra $\si{\metre\per\second\squared}$ til $\si{\kilo\metre\per\minute\squared}$. Det er lurt å regne om lengdeenhet og tidsenhet hver for seg (se \textit{Å regne med enheter} i Kapittel~\ref{komp-Chapter:hvaerfysikk}).

Lengde: $\SI{1}{\metre} = 10^{-3}\,\si{\kilo\metre}$.

Tid: $\SI{1}{\second} = \tfrac{1}{60}\,\si{\minute}$, så $\SI{1}{\second\squared} = \tfrac{1}{3600}\,\si{\minute\squared}$, altså $\dfrac{1}{\si{\second\squared}} = 3600\,\dfrac{1}{\si{\minute\squared}}$.

Vi setter inn begge omregningene:
\[
    \mathrm{g}_\textrm{M} = \SI{1.62}{\metre\per\second\squared} = 1.62\cdot 10^{-3}\,\si{\kilo\metre}\cdot 3600\,\si{\per\minute\squared} = \doubleunderline{\SI{5.83}{\kilo\metre\per\minute\squared}}
\]
}

\oppgaver{Signifikante siffer}{signsiff}{%
%
\begin{enumerate}[a)]
    \item Oppgi antall signifikante siffer i: $0.00420$, \ $3.14159$, \ $1\,200$ og $165.00$.
    \item Regn ut $12.3 \cdot 4.56$ og oppgi svaret med riktig antall signifikante siffer.
    \item Regn ut $18.0 / 4.5$ og oppgi svaret med riktig antall signifikante siffer.
\end{enumerate}
%
}

\svar[title={Løsningsforslag}]{%
\noindent
\textbf{a)} Antall signifikante siffer (se \textit{Signifikante siffer} i Kapittel~\ref{komp-Chapter:hvaerfysikk}):
\begin{itemize}
    \item $0.00420$: de ledende nullene teller ikke, men den etterfølgende nullen gjør det $\Rightarrow$ \doubleunderline{3 signifikante siffer}.
    \item $3.14159$: alle sifre er signifikante $\Rightarrow$ \doubleunderline{6 signifikante siffer}.
    \item $1\,200$: uten desimaltegn regnes de to siste nullene som ikke-signifikante $\Rightarrow$ \doubleunderline{2 signifikante siffer}.
    \item $165.00$: desimaltegnet gjør at alle nuller teller $\Rightarrow$ \doubleunderline{5 signifikante siffer}.
\end{itemize}

\smallskip\noindent
\textbf{b)} Begge faktorene har 3 signifikante siffer, så svaret skal også ha 3:
\[
    12.3 \cdot 4.56 = 56.088 \approx \doubleunderline{56.1}
\]

\smallskip\noindent
\textbf{c)} $18.0$ har 3 signifikante siffer, $4.5$ har 2, så svaret skal ha 2:
\[
    18.0 / 4.5 = \doubleunderline{4.0}
\]
}

\oppgaver{Standardform}{standardform-opg}{%
%
Skriv følgende tall på standardform:
\begin{enumerate}[a)]
    \item $0.000\,045$
    \item $3\,750\,000$
    \item $0.100$
    \item $299\,792\,458$
\end{enumerate}
Oppgi også antall signifikante siffer i hvert tall, slik du tolker det.
%
}

\svar[title={Løsningsforslag}]{%
Et tall på standardform skrives $T = t\cdot 10^n$ med $1\le t < 10$, se \textit{Tall på standardform} i kompendiets Kapittel~\ref{komp-Chapter:hvaerfysikk}.
\begin{enumerate}[a)]
    \item $0.000\,045 = \doubleunderline{4.5\cdot 10^{-5}}$ (2 signifikante siffer)
    \item $3\,750\,000 = \doubleunderline{3.75\cdot 10^{6}}$ (3 signifikante siffer)
    \item $0.100 = \doubleunderline{1.00\cdot 10^{-1}}$ (3 signifikante siffer, desimaltegnet gjør nullene signifikante)
    \item $299\,792\,458 = \doubleunderline{2.99792458\cdot 10^{8}}$ (alle 9 sifrene er signifikante --- dette er for øvrig definisjonen av lyshastigheten $c$)
\end{enumerate}
}

\oppgaver{Dekadiske prefikser}{dekpref-opg}{%
%
\begin{enumerate}[a)]
    \item Et batteri har kapasitet $\SI{4\,200}{\milli\ampere\hour}$ (milliamperetimer). Skriv dette som et tall i Ah og på standardform.
    \item En avstand måles til $\SI{0.043}{\milli\metre}$. Skriv dette i $\mu\textrm{m}$ (mikrometer) og i nm (nanometer).
    \item Lyset bruker ca. $\SI{499}{\second}$ fra solen til jorda. Lyset har fart $c = 3.00 \cdot 10^8\,\si{\metre\per\second}$. Finn avstanden fra solen til jorda i km og skriv svaret på standardform.
\end{enumerate}
%
}

\svar[title={Løsningsforslag}]{%
\noindent
\textbf{a)} $\SI{4200}{\milli\ampere\hour} = 4200\cdot 10^{-3}\,\si{\ampere\hour} = \SI{4.2}{\ampere\hour} = \doubleunderline{4.2\cdot 10^{0}\,\si{\ampere\hour}}$

\smallskip\noindent
\textbf{b)} Vi bruker at $\SI{1}{\milli\metre} = 10^3\,\si{\micro\metre} = 10^6\,\si{\nano\metre}$:
\[
    \SI{0.043}{\milli\metre} = 0.043\cdot 10^3\,\si{\micro\metre} = \doubleunderline{\SI{43}{\micro\metre}} = 0.043\cdot 10^6\,\si{\nano\metre} = \doubleunderline{\SI{4.3e4}{\nano\metre}}
\]

\smallskip\noindent
\textbf{c)} Avstanden er fart multiplisert med tid:
\[
    d = c\cdot t = 3.00\cdot 10^8\,\si{\metre\per\second}\cdot \SI{499}{\second} = 1.497\cdot 10^{11}\,\si{\metre} = 1.497\cdot 10^{8}\,\si{\kilo\metre} \approx \doubleunderline{1.50\cdot 10^{8}\,\si{\kilo\metre}}
\]
(avrundet til 3 signifikante siffer, siden $c$ og $t$ begge er oppgitt med 3 signifikante siffer).
}

\oppgaver{Fysisk størrelse}{k1oppg7}{%
%
Hva skal stå på venstre side av likhetstegnet?
\[
    \rule{4cm}{0.4pt} \;=\; \textrm{verdi} \cdot \textrm{enhet}
\]
%
}

\svar[title={Løsningsforslag}]{%
På venstre side skal det stå \doubleunderline{\textrm{Fysisk størrelse}}, slik at likningen blir
\[
    \textrm{Fysisk størrelse} = \textrm{verdi}\cdot\textrm{enhet},
\]
som er selve definisjonen av en fysisk størrelse (se \textit{Fysiske størrelser} i Kapittel~\ref{komp-Chapter:hvaerfysikk}).
}

\oppgaver{Dekadiske prefikser -- matching}{k1oppg8}{%
%
Fyll inn faktoren for hvert prefiks i tabellen:
\begin{center}
\renewcommand{\arraystretch}{1.3}
\begin{tabular}{|l|c|c|}
\hline
\rowcolor{blue!20!black!10}
\textbf{Prefiks} & \textbf{Symbol} & \textbf{Faktor} \\
\hline
mega  & M & \\[4pt]
nano  & n & \\[4pt]
milli & m & \\[4pt]
centi & c & \\[4pt]
hekto & h & \\[4pt]
kilo  & k & \\[4pt]
giga  & G & \\[4pt]
\hline
\end{tabular}
\end{center}
%
}

\svar[title={Løsningsforslag}]{%
Faktorene er gitt i Tabell~\ref{komp-tab:dekPref} i kompendiet:
\begin{center}
\renewcommand{\arraystretch}{1.3}
\begin{tabular}{|l|c|c|}
\hline
\rowcolor{blue!20!black!10}
\textbf{Prefiks} & \textbf{Symbol} & \textbf{Faktor} \\
\hline
mega  & M & $10^{6}$ \\[2pt]
nano  & n & $10^{-9}$ \\[2pt]
milli & m & $10^{-3}$ \\[2pt]
centi & c & $10^{-2}$ \\[2pt]
hekto & h & $10^{2}$ \\[2pt]
kilo  & k & $10^{3}$ \\[2pt]
giga  & G & $10^{9}$ \\[2pt]
\hline
\end{tabular}
\end{center}
}

\oppgaver{Standardform}{k1oppg9}{%
%
Hvordan ser tallet $100.1$ ut når det er skrevet på standardform?
%
}

\svar[title={Løsningsforslag}]{%
Vi flytter desimaltegnet slik at det blir ett siffer foran det, og teller hvor mange plasser vi flyttet:
\[
    100.1 = \doubleunderline{1.001\cdot 10^{2}}
\]
}

\oppgaver{Signifikante siffer}{k1oppg10}{%
%
Per kjøper 1523 skruer til garasjen sin. Hver skrue veier $\SI{10.1}{\gram}$.
Hvor mange signifikante sifre kan vi angi totalvekten til skruene med?
%
}

\svar[title={Løsningsforslag}]{%
Antallet skruer, $1523$, er et \textit{opptalt} antall og er dermed eksakt (uendelig mange signifikante siffer) --- det er ikke dette tallet som begrenser nøyaktigheten. Vekten per skrue, $\SI{10.1}{\gram}$, er derimot en \textit{målt} størrelse med 3 signifikante siffer.

Totalvekten er
\[
    m_\textrm{tot} = 1523\cdot\SI{10.1}{\gram} = \SI{15382.3}{\gram},
\]
men siden den målte størrelsen bare har 3 signifikante siffer, kan ikke svaret oppgis med flere enn \doubleunderline{3 signifikante siffer}:
\[
    m_\textrm{tot} \approx \SI{1.54e4}{\gram} = \SI{15.4}{\kilo\gram}.
\]
}
